\section*{Abstract}
\addcontentsline{toc}{section}{Abstract}

Accurate three-axis attitude determination and control is a foundational requirement for nearly all spacecraft missions, and it becomes especially demanding for CubeSats, whose small size, limited actuation authority, and tight power budgets constrain the achievable pointing performance. This report presents the development and analysis of a quaternion-based attitude dynamics and control simulation for a 3U CubeSat equipped with a three-axis reaction wheel assembly. The spacecraft attitude kinematics are propagated using unit quaternions to avoid the singularities inherent to three-parameter Euler angle representations, while the rotational dynamics are modeled using Euler's rigid-body equations of motion about the spacecraft's principal axes. A proportional-derivative (PD) control law is formulated directly on the quaternion attitude error and body angular velocity to command reaction wheel torques, and is shown analytically, via a Lyapunov argument, to be almost-globally asymptotically stable. The complete model is implemented in MATLAB using a fourth-order Runge-Kutta numerical integrator and exercised over a representative post-deployment scenario combining an initial $19.7^\circ$/s tumble with a $120^\circ$ large-angle slew to a target inertial attitude. Starting from this combined initial condition, the closed-loop simulation converges to within $0.11^\circ$ of the target attitude and $0.01^\circ$/s of zero angular rate within \SI{74.1}{\second}, with the reaction wheels briefly saturating at their \SI{1}{\milli\newton\meter} torque limit for the first \SI{5.6}{\second}. The results are validated against three independent checks: conservation of the quaternion unit-norm constraint (drift held at floating-point round-off level), agreement between the primary fixed-step integrator and an independent adaptive solver (maximum discrepancy $1.3\times10^{-4}$ degrees), and monotonic decrease of the derived Lyapunov function, confirming the simulated closed-loop behavior matches its theoretical stability certificate. Simplifying assumptions -- including a rigid-body spacecraft, idealized torque-source actuators, and the exclusion of environmental disturbance torques -- are stated explicitly, and their impact on the fidelity of the results is discussed.

\clearpage
